\section{Methods}
\subsection{Encoder and Tuple Objective}
For nonnegative convolutional activations $x_{c,u}$ at channel $c$ and position $u\in\mathcal U$, GeM computes
\begin{equation}
 g_c=\left(\frac{1}{|\mathcal U|}\sum_{u\in\mathcal U}x_{c,u}^{p}\right)^{1/p},
 \qquad \mathbf z=\frac{\mathbf g}{\|\mathbf g\|_2}.
\end{equation}
A shared learnable $p$ starts at 3. VGG16 features, without the classifier and final pooling, produce a 512-dimensional normalized descriptor. The initial ranking uses cosine score $s_i=\mathbf z_q^\top\mathbf z_i$.

Each training tuple contains a query $q$, one metadata-defined positive $p$, and five negatives $n_j$. The implemented objective is
\begin{equation}
 \mathcal L=\tfrac12 d(q,p)^2+\tfrac12\sum_{j=1}^{5}[m-d(q,n_j)]_+^2,
 \quad m=0.7,
\end{equation}
where $[a]_+=\max(a,0)$ and $d$ is Euclidean descriptor distance. The implementation adds $10^{-6}$ to coordinate differences before computing the norm for numerical stabilization. Positive relations come from explicit metadata pairs; cluster membership alone does not assign a positive label. Hard-negative mining uses current clean descriptors, excludes the query cluster, and selects negatives from distinct clusters.

\subsection{Training Modifications}
\Rone{} adds independent photometric and degradation transforms to the images used in the training forward pass. With probability 0.8, brightness and contrast factors are drawn from $[0.8,1.2]$ and color from $[0.85,1.15]$, in random order. Gaussian blur is applied with probability 0.15 and radius $[0.2,1.0]$ pixels; JPEG encoding with probability 0.2 and integer quality 55--95; and erasing with probability 0.15, area 1--5\%, and aspect ratio 0.5--2. Erasing fills zero after normalization, corresponding to the normalization mean in RGB. Clean mining and validation do not use these transforms.

\Rtwo{} adds a spatial gate before GeM. For feature tensor $X$, the gate is
\begin{equation}
 A(X)=2\sigma\!\left(\operatorname{Conv}_{7\times7}
 [\operatorname{mean}_cX;\operatorname{max}_cX]\right).
\end{equation}
The two-input, one-output convolution has 99 parameters including bias. Zero-initialized weights and bias give $A(X)=1$, so the initial gated tensor equals $X$. The output descriptor dimensionality remains unchanged.

\subsection{Query Expansion and Feedback Selection}
Let $\mathcal F$ be positive-cosine candidates from the first top-$k$ ranking. We compute
\begin{equation}
 \mathbf q'=\operatorname{norm}_2\!\left(\mathbf q+
 \lambda\frac{\sum_{i\in\mathcal F}w_i\mathbf z_i}{\sum_{i\in\mathcal F}w_i}\right),
 \quad w_i=\max(s_i,0)^\alpha .
 \label{eq:qe}
\end{equation}
AQE uses $\alpha=0$; weighted QE uses $\alpha=3$. We fix $k=10$ and $\lambda=0.5$ for one feedback iteration. Gated QE additionally requires $s_i\geq0.5$ and at least three accepted candidates, otherwise returning the original query. Setting $\lambda=0$ also returns the baseline. The residual form in~\eqref{eq:qe} is stated explicitly because conventions for mixing query and feedback differ.

\subsection{Local and Geometric Evidence}
We extract up to 2,000 RootSIFT keypoints per image. Matching uses a 0.75 nearest-neighbor ratio test and mutual nearest-neighbor consistency. Homography RANSAC uses a 3-pixel threshold, confidence 0.995, and at most 2,000 iterations. A valid result requires at least eight inliers and a finite, full-rank homography with condition number below $10^{12}$.

For a top-50 candidate with $M_i$ matches and $I_i$ inliers, local re-ranking uses evidence $e_i=\min(M_i/50,1)$. Geometric re-ranking instead uses
\begin{equation}
 e_i=\frac{I_i}{\max(M_i,1)}\min(I_i/20,1)
\end{equation}
when geometry is valid, and zero otherwise. The combined score is $(1-\eta)(s_i+1)/2+\eta e_i$, with $\eta=0.5$. Only the candidate prefix changes; the remaining ranking is retained. Invalid query features or no valid geometric candidate trigger a logged fallback where applicable.

\emph{Verified QE} checks the first 50 global candidates, selects at most ten with valid geometry and positive cosine, and applies~\eqref{eq:qe}. Fewer than three valid feedback images returns the original query. Search is then repeated over the entire gallery, without a second re-ranking stage. Relevance and junk labels never select feedback.

\subsection{Fusion and Diffusion}
Score fusion averages the two encoders' cosine scores after image-ID alignment. RRF uses both full rankings:
\begin{equation}
 S_{\mathrm{RRF}}(i)=\sum_{r=1}^{2}\frac{1}{60+\operatorname{rank}_r(i)}.
\end{equation}
We do not average raw embeddings across encoders.

Each encoder also defines a gallery graph with $k=20$, no self-edges, and directed weights $\max(s_{ij},0)^3$. Union and mutual graphs use the elementwise maximum and minimum with the transposed affinity, respectively. Let $S=D^{-1/2}WD^{-1/2}$, using zero inverse degree for isolated nodes. Initial query scores $y_i$ are positive-cosine cubes normalized to unit sum. We iterate
\begin{equation}
 \mathbf f^{(t+1)}=(1-\beta)\mathbf y+\beta S\mathbf f^{(t)},
 \quad \beta=0.9,
\end{equation}
starting at $\mathbf f^{(0)}=\mathbf y$, stopping at relative L2 change below $10^{-6}$ or 100 iterations. Non-finite results or absent initial mass fall back to cosine. Labels are not used in graph construction or propagation.
